\documentclass{beamer}
\usepackage[utf8]{inputenc}
\usetheme{madrid}

\usepackage{mathtools}
\usepackage{amssymb,amsmath,amsthm}
\usepackage[mathscr]{eucal}
\usepackage[procnames]{listings}
\usepackage{color}
\usepackage{pdfpages}
\usepackage{pgfplots}
\usepackage{pgfplotstable} 
\pgfplotsset{compat=1.18}
\usepackage{tikz}
\usepackage{caption}
\usepackage{extarrows}
\usepackage{units}
\usepackage{proof}
\usepackage{algorithm}
\usepackage{algpseudocode}
\usetikzlibrary{positioning}

\title{Second talk: Acyclic Graph coloring via SAT}
\author{Nora Depenheuer}
\date{6.8.26}

\begin{document}
	\begin{frame}
		\titlepage
	\end{frame}
	\begin{frame}
		\frametitle{Overview}
		\begin{enumerate}
			\item Problem
			\item NP-hardness
			\item The general algorithm
			\item Lower and upper bound
			\item Preprocessing
			\item Specific SAT formulations
			\item Variants and their runtime
		\end{enumerate}
	\end{frame}
	\begin{frame}
		\begin{definition}[Acyclic $k$-coloring]
			Given graph $G=(V,E)$ an \underline{acyclic $k$-coloring} is a $k$-coloring c (i.e. a map $c:V\to \{1,...k\}$ where $\forall \{v,w\}\in E: c(v)\neq c(w)$) with no bicolored cycle.
		\end{definition}
		\begin{definition}
			The \underline{acyclic coloring number} is the minimum $k$ such that there is an acyclic $k$-coloring.
		\end{definition}
	\end{frame}
	\begin{frame}
		To justify, why we consider solving this via SAT, we first prove, that our problem is NP-hard. For this we prove that:
		\begin{theorem}
			Deciding whether there is an acyclic $k$-coloring for given $k\geq 3$ is NP-complete. 
		\end{theorem}  
		\begin{proof}
			\only<1>{Seeing that this problem is in NP is easy, since checking for whether something is a coloring is possible in $O(m)$ and checking for bicolored cycles is possible in $O(k^2(n+m))$. \\
			To prove NP-completeness, we reduce an instance of regular decisional coloring to the acyclic one. So consider a Graph $G=(V,E)$ and a number $k$. Construct a new graph $G'$ by adding for each edge $\{v,w\}\in E$ the vertices $\{v,w,u_1,...u_k\}$ to $V'$ and edges $\{v,u_i\},\{w,u_i\}\forall 1\leq i\leq k$ to $E'$.}
			\only<2>{Now prove that if there is a $k$-coloring $c$ on $G$, then there is an acyclic $k$-coloring $c'$ on $G'$. For this, assume a $c$ and construct $c'$ by setting $c'(v)=c(v)\forall v\in V$. For the $u_i$ created from some edge $\{v,w\}\in E$ set $c'(v)\neq c'(u_i)\neq c'(w)$ (obviously possible since $k\geq 3$). Then $c'$ is clearly a $k$-coloring and since any cycle must contain a path of form $(v,u_i,w)$, it must contain at least 3 colors.}
			\only<3>{Now prove that if there is an acyclic $k$-coloring $c'$ on $G'$, then there is a $k$-coloring $c$ on $G$. For this, assume a $c'$ and construct $c$ by setting $c(v)=c'(v)\forall v\in V$. Now suppose there is an edge $\{v,w\}\in E: c(v)=c(w)$ then all $u_i$ created from edge $\{v,w\}$ must have different colors in $c'$ since, if there are $u_i,u_j$ with $c'(u_i)=c'(u-j)$ then $v,u_i,w,u_j$ is a bicolored cycle. This is not possible. If all $u_i$ use different colors, then $c'$ uses $k+1$ colors which is a contradiction.}
			\alt<3>{\qedhere}{\phantom\qedhere}
		\end{proof}
	\end{frame}
	
	\begin{frame}
		\frametitle{The general algorithm}
		\begin{algorithm}[H]
			\caption{General approach for Graph $G$} 
			\begin{algorithmic}[1]
				\State $lb =$ find lower bound $(G)$
				\State preprocessing $(G, lb)$
				\State $ub=$ find upper bound $(G)$
				\For {$k=ub-1,\ldots,lb$}
				\State Initialize SAT $(G,k)$
				\State sat $=$ Solve SAT
				\While{sat $=$ true and $\exists$bicolored cycle}
				\State Update SAT $(G,k,cycle)$
				\State sat $=$ Solve SAT
				\EndWhile
				\If {sat $=$ false}
				\State return $k-1$
				\EndIf
				\EndFor
				\State return $lb$
			\end{algorithmic} 
		\end{algorithm}	
		We generally find cycles by running BFS on all pairs of colors.
	\end{frame}
	\begin{frame}
		\frametitle{Lower and upper bound}
		We know from regular coloring, that the clique number is a viable lower bound. Since we only have more restraints, this still holds. We compute the clique number via CliSAT.\\~\\
		For the upper bound we use acyclic DSatur as presented in a previous Master's thesis. This means we simply go through all vertices in a certain order and assign them the next possible color, creating a new one whenever needed.
	\end{frame}
	\begin{frame}
		\frametitle{Preprocessing}
		We have two rules which we apply until no further changes happen:
		\begin{lemma}
			We can remove vertices whose neighborhood form a clique.
		\end{lemma}
		\begin{lemma}
			We can remove vertices whose second neighborhood ($N_2(v)=\{u: dist(v,u)\leq 2\}$) has at most lower bound many vertices and $|N(v)|<lb$.
		\end{lemma}
	\end{frame}
	\begin{frame}
		\begin{lemma}
			We can remove vertices whose neighborhood form a clique.
		\end{lemma}
		\begin{proof}
			Assume we have an acyclic coloring on $G-v$.
			$N(v)$ is a clique so in particular, every vertex in it will be assigned a different color. Thus by coloring $v$ in a color not yet represented in $N(v)$, we get a valid coloring. Since any cycle going through $v$ has to pass through at least two of it's neighbors, the cycle already contains 3 colors and is thus not bicolored.\\
			If $|N(v)|<lb$ then there is one more color with which to color $v$. $|N(v)|=lb$ is not possible, since $N(v)\cup\{v\}$ is a clique and we compute the lower bound to be the maximum clique size.
		\end{proof}
	\end{frame}
	\begin{frame}
		\begin{lemma}
			We can remove vertices whose second neighborhood ($N_2(v)=\{u: dist(v,u)\leq 2\}\setminus \{v\}$) has at most lower bound many vertices and $|N(v)|<lb$.
		\end{lemma}
		\begin{proof}
			Assume we have an acyclic coloring on $G-v$.
			W.l.o.g $|N_2(v)|=k$. If it's lower, then there is a color that is not used in $N_2(v)$ and thus by assigning it to v, we get a viable acyclic coloring. Differentiate two cases:
			\begin{enumerate}
				\item If all vertices in $N_2(v)$ get a different color, then there must be a vertex in $N_2(v)\setminus N(v)$ whose color we can assign to $v$. This is obviously a viable coloring and it is acyclic since any cycles through $v$ has at least 3 colors.
				\item There are some vertices in $N_2(v)$ with the same color. Then the amount of colors used in $N_2(v)$ is less than $lb$ and there is one color not represented that we can assign to $v$.
			\end{enumerate}
		\end{proof}
	\end{frame}
	\begin{frame}
		\frametitle{Assignment formulation}
		Have $n\cdot k$ variables. With $x_{v,i}=\begin{cases*}$1$ & if $v$ has the color $i$\\ 0 & else\end{cases*}$. 
		
		\begin{alignat*}{2}
			&\sum_{i=0}^{k}x_{v,i}=1 &\forall v\in V\\
			&x_{v,i}+x_{w,i}\leq 1 &\forall \{v,w\}\in E,\forall i\in \{1,...,k\}
		\end{alignat*}\\
		
		As logic formula this corresponds to the conjunction of the following clauses
		
		\begin{alignat*}{2}
			&\bigvee_{i=0}^{k}x_{v,i} &\forall v\in V\\
			&\neg x_{v,i}\vee \neg x_{v,j} &\forall v\in V,\forall i<j\in\{1,...,k\}\\
			&\neg x_{v,i}\vee\neg x_{w,i} &\forall \{v,w\}\in E,\forall i\in \{1,...,k\}
		\end{alignat*}\\
		
		To forbid a cycle $C=(c_1,...c_l)\subseteq V$, we then add the clauses
		\begin{alignat*}{2}
			&\bigvee_{s=1,3,5,..l-1}\neg x_{c_s,i}\vee\bigvee_{s=2,4,6,..l}\neg x_{c_s,j} &\forall i\neq j\in\{1,...,k\}
		\end{alignat*}
	\end{frame}
	\begin{frame}
		\frametitle{Assignment formulation}
		Have $n\cdot k$ variables. With $x_{v,i}=\begin{cases*}$1$ & if $v$ has the color $i$\\ 0 & else\end{cases*}$. 		
		\begin{alignat*}{2}
			&\bigvee_{i=0}^{k}x_{v,i} &\forall v\in V\\
			&\neg x_{v,i}\vee \neg x_{v,j} &\forall v\in V,\forall i<j\in\{1,...,k\}\\
			&\neg x_{v,i}\vee\neg x_{w,i} &\forall \{v,w\}\in E,\forall i\in \{1,...,k\}
		\end{alignat*}\\
		
		To forbid a cycle $C=(c_1,...c_l)\subseteq V$, we then add the clauses
		\begin{alignat*}{2}
			&\bigvee_{s=1,3,5,..l-1}\neg x_{c_s,i}\vee\bigvee_{s=2,4,6,..l}\neg x_{c_s,j} &\forall i\neq j\in\{1,...,k\}
		\end{alignat*}
	\end{frame}
	\begin{frame}
		\frametitle{Partial-ordering formulation}
		Have $n\cdot k$ variables. We assume an ordering on the colors (by naming them in $\{1,...,k\}$) and then $x_{v,i}=\begin{cases*}$1$ & if $v$ has a color bigger than $i$\\ 0 & if $v$ has a color smaller or equal to $i$\end{cases*}$. So every vertex will be assigned the color with $x_{v,i-1}-x_{v,i}=1$. 		
		\begin{alignat*}{2}
			&x_{v,k}=0 &\forall v\in V\\
			&x_{v,i}-x_{v,i+1}\geq 0 &\forall v\in V,\forall i\in\{1,...,k-1\}\\
			&x_{v,1}+x_{w,1}\geq 1 &\forall \{v,w\}\in E\\
			&x_{v,i-1}-x_{v,i}+x_{w,i-1}-x_{i}\leq 1 &\forall \{v,w\}\in E,\forall i\in \{2,...,k\}
		\end{alignat*}\\		
		As logic formula this corresponds to the conjunction of the following clauses		
		\begin{alignat*}{2}
			&\neg x_{v,k} &\forall v\in V\\
			&x_{v,i}\vee\neg x_{v,i+1} &\forall v\in V,\forall i\in\{1,...,k-1\}\\
			&x_{v,1}\vee x_{w,1} &\forall \{v,w\}\in E\\
			&\neg x_{v,i-1}\vee x_{v,i}\vee\neg x_{w,i-1}\vee x_{i} &\forall \{v,w\}\in E,\forall i\in \{2,...,k\}
		\end{alignat*}
	\end{frame}
	\begin{frame}
		\frametitle{Partial-ordering formulation}
		Have $n\cdot k$ variables. With $x_{v,i}=\begin{cases*}$1$ & if $v$ has a color bigger than $i$\\ 0 & if $v$ has a color smaller or equal to $i$\end{cases*}$.\\ So every vertex will be assigned the color with $x_{v,i-1}-x_{v,i}=1$. 		
		\begin{alignat*}{2}
			&\neg x_{v,k} &\forall v\in V\\
			&x_{v,i}\vee\neg x_{v,i+1} &\forall v\in V,\forall i\in\{1,...,k-1\}\\
			&x_{v,1}\vee x_{w,1} &\forall \{v,w\}\in E\\
			&\neg x_{v,i-1}\vee x_{v,i}\vee\neg x_{w,i-1}\vee x_{i} &\forall \{v,w\}\in E,\forall i\in \{2,...,k\}
		\end{alignat*}\\		
		To forbid a cycle $C=(c_1,...c_l)\subseteq V$, we then add the clause
		\begin{alignat*}{2}
			&\bigvee_{s=1,3,5,..l-1}\neg x_{c_s,i-1}\vee x_{c_s,i}\vee\bigvee_{s=2,4,6,..l}\neg x_{c_s,j-1}\vee x_{c_s,j} &i=c(c_1),j=c(c_2)
		\end{alignat*}
		though, if one color w.l.o.g $i$ is $1$, there is no $x_{c_s,i-1}$ so leave that part out and just require $x_{c_s,i}$.
	\end{frame}
	\begin{frame}
		\frametitle{Runtime}
		Letting both of these run on Erdös-Rényi graphs with $n=20$ and varying density $d=0.3$ and $d=0.5$ we get the median:
		\begin{figure}[H]
		\centering
		\begin{minipage}{.5\textwidth}
			\centering
			\pgfplotstableread[col sep=semicolon]{Evaluating data 1.csv}\rawdata
			\resizebox{\columnwidth}{!}{
			\begin{tikzpicture}
				\begin{axis} [
					ybar,
					bar width=20pt,
					ylabel={Time (ms)},
					xlabel={density $0.3$},
					% Automatically read category text labels directly from the table
					xtick=data,
					xticklabels from table={\rawdata}{category},
					ymin=0, ymax=2500,
					enlarge x limits=0.3
					]
					
					% 2. Clean plot generation using structured table columns
					\addplot+[
					error bars/.cd,
					y dir=both,
					y explicit
					] table [
					x expr=\coordindex, 
					y=median2,
					% Dynamically calculate lower and upper error whisker margins
					y error minus expr=\thisrow{median2}-\thisrow{min},
					y error plus expr=\thisrow{max}-\thisrow{median2}
					] {\rawdata};
					
				\end{axis}
			\end{tikzpicture}}
		\end{minipage}%
		\begin{minipage}{.5\textwidth}
			\centering
			\pgfplotstableread[col sep=semicolon]{Evaluating data 2.csv}\rawdata
			\resizebox{\columnwidth}{!}{
			\begin{tikzpicture}
				\begin{axis} [
					ybar,
					bar width=20pt,
					ylabel={Time (ms)},
					xlabel={density $0.5$},
					% Automatically read category text labels directly from the table
					xtick=data,
					xticklabels from table={\rawdata}{category},
					ymin=0, ymax=2500000,
					enlarge x limits=0.3
					]
					
					% 2. Clean plot generation using structured table columns
					\addplot+[
					error bars/.cd,
					y dir=both,
					y explicit
					] table [
					x expr=\coordindex, 
					y=median2,
					% Dynamically calculate lower and upper error whisker margins
					y error minus expr=\thisrow{median2}-\thisrow{min},
					y error plus expr=\thisrow{max}-\thisrow{median2}
					] {\rawdata};
					
				\end{axis}
			\end{tikzpicture}}
		\end{minipage}
	\end{figure}
	for $d=0.8$ the assignment version does not run through in any acceptable runtime.
	\end{frame}
	\begin{frame}
		\frametitle{Runtime}
		Letting both of these run on Erdös-Rényi graphs with $n=20$ and varying density $d=0.3$ and $d=0.5$ we get the median:
		\begin{figure}[H]
			\centering
			\begin{minipage}{.5\textwidth}
				\centering
				\pgfplotstableread[col sep=semicolon]{Evaluating data 1.csv}\rawdata
				\resizebox{\columnwidth}{!}{
					\begin{tikzpicture}
						\begin{axis} [
							ybar,
							bar width=20pt,
							ylabel={Time (ms)},
							xlabel={density $0.3$},
							% Automatically read category text labels directly from the table
							xtick=data,
							xticklabels from table={\rawdata}{category},
							ymin=0, ymax=40,
							enlarge x limits=0.3
							]
							
							% 2. Clean plot generation using structured table columns
							\addplot+[
							error bars/.cd,
							y dir=both,
							y explicit
							] table [
							x expr=\coordindex, 
							y=median2] {\rawdata};
							
						\end{axis}
				\end{tikzpicture}}
			\end{minipage}%
			\begin{minipage}{.5\textwidth}
				\centering
				\pgfplotstableread[col sep=semicolon]{Evaluating data 2.csv}\rawdata
				\resizebox{\columnwidth}{!}{
					\begin{tikzpicture}
						\begin{axis} [
							ybar,
							bar width=20pt,
							ylabel={Time (ms)},
							xlabel={density $0.5$},
							% Automatically read category text labels directly from the table
							xtick=data,
							xticklabels from table={\rawdata}{category},
							ymin=0, ymax=50000,
							enlarge x limits=0.3
							]
							
							% 2. Clean plot generation using structured table columns
							\addplot+[
							error bars/.cd,
							y dir=both,
							y explicit
							] table [
							x expr=\coordindex, 
							y=median2] {\rawdata};
							
						\end{axis}
				\end{tikzpicture}}
			\end{minipage}
		\end{figure}
		for $d=0.8$ the assignment version does not run through in any acceptable runtime. So we continue from here with the partial ordering and try to improve it.
	\end{frame}
	\begin{frame}
		\frametitle{Variants}
		So there were several incremental small changes made to the program to improve the runtime.
		\begin{enumerate}
			\item finding more than one cycle per iteration
			\item when finding cycles first ordering the pairs of colors by how many edges with these colors there are
			\item remembering all cycles found in a previous iteration with different $k$ and initially forbidding them
		\end{enumerate}
	\end{frame}
	\begin{frame}
		\frametitle{Runtime}
		Letting all of these run on Erdös-Rényi graphs with $n=20$ and varying density $d=0.3$ and $d=0.5$ we get the median:
		\begin{figure}[H]
			\centering
			\begin{minipage}{.5\textwidth}
				\centering
				\pgfplotstableread[col sep=semicolon]{Evaluating data 3.csv}\rawdata
				\resizebox{\columnwidth}{!}{
					\begin{tikzpicture}
						\begin{axis} [
							ybar,
							bar width=20pt,
							ylabel={Time (ms)},
							xlabel={density $0.3$},
							% Automatically read category text labels directly from the table
							xtick=data,
							xticklabels from table={\rawdata}{category},        xticklabel style={rotate=60, anchor=north east}],
							ymin=0, ymax=2500,
							enlarge x limits=0.3
							]
							
							% 2. Clean plot generation using structured table columns
							\addplot+[
							error bars/.cd,
							y dir=both,
							y explicit
							] table [
							x expr=\coordindex, 
							y=median2,
							% Dynamically calculate lower and upper error whisker margins
							y error minus expr=\thisrow{median2}-\thisrow{min},
							y error plus expr=\thisrow{max}-\thisrow{median2}
							] {\rawdata};
							
						\end{axis}
				\end{tikzpicture}}
			\end{minipage}%
			\begin{minipage}{.5\textwidth}
				\centering
				\pgfplotstableread[col sep=semicolon]{Evaluating data 4.csv}\rawdata
				\resizebox{\columnwidth}{!}{
					\begin{tikzpicture}
						\begin{axis} [
							ybar,
							bar width=20pt,
							ylabel={Time (ms)},
							xlabel={density $0.5$},
							% Automatically read category text labels directly from the table
							xtick=data,
							xticklabels from table={\rawdata}{category},        xticklabel style={rotate=60, anchor=north east}],
							ymin=0, ymax=250000,
							enlarge x limits=0.3
							]
							
							% 2. Clean plot generation using structured table columns
							\addplot+[
							error bars/.cd,
							y dir=both,
							y explicit
							] table [
							x expr=\coordindex, 
							y=median2,
							% Dynamically calculate lower and upper error whisker margins
							y error minus expr=\thisrow{median2}-\thisrow{min},
							y error plus expr=\thisrow{max}-\thisrow{median2}
							] {\rawdata};
							
						\end{axis}
				\end{tikzpicture}}
			\end{minipage}
		\end{figure}
	\end{frame}
	\begin{frame}
		\frametitle{Runtime}
		Letting all of these run on Erdös-Rényi graphs with $n=20$ and varying density $d=0.3$ and $d=0.5$ we get the median:
		\begin{figure}[H]
			\centering
			\begin{minipage}{.5\textwidth}
				\centering
				\pgfplotstableread[col sep=semicolon]{Evaluating data 3.csv}\rawdata
				\resizebox{\columnwidth}{!}{
					\begin{tikzpicture}
						\begin{axis} [
							ybar,
							bar width=20pt,
							ylabel={Time (ms)},
							xlabel={density $0.3$},
							% Automatically read category text labels directly from the table
							xtick=data,
							xticklabels from table={\rawdata}{category},        xticklabel style={rotate=60, anchor=north east}],
							ymin=0, ymax=40,
							enlarge x limits=0.3
							]
							
							% 2. Clean plot generation using structured table columns
							\addplot+[
							error bars/.cd,
							y dir=both,
							y explicit
							] table [
							x expr=\coordindex, 
							y=median2]{\rawdata};
							
						\end{axis}
				\end{tikzpicture}}
			\end{minipage}%
			\begin{minipage}{.5\textwidth}
				\centering
				\pgfplotstableread[col sep=semicolon]{Evaluating data 4.csv}\rawdata
				\resizebox{\columnwidth}{!}{
					\begin{tikzpicture}
						\begin{axis} [
							ybar,
							bar width=20pt,
							ylabel={Time (ms)},
							xlabel={density $0.5$},
							% Automatically read category text labels directly from the table
							xtick=data,
							xticklabels from table={\rawdata}{category},        xticklabel style={rotate=60, anchor=north east}],
							ymin=0, ymax=25000,
							enlarge x limits=0.3
							]
							
							% 2. Clean plot generation using structured table columns
							\addplot+[
							error bars/.cd,
							y dir=both,
							y explicit
							] table [
							x expr=\coordindex, 
							y=median2] {\rawdata};
							
						\end{axis}
				\end{tikzpicture}}
			\end{minipage}
		\end{figure}
		These still aren't very fast and for $n=30$ take hours. So we try something different.
	\end{frame}
	\begin{frame}
		\frametitle{Representative formulation}
		Have three types of variables. First $n$ many $x_u$, where $x_u=\begin{cases*}$1$ & if $v$ is a representative\\ 0 & else\end{cases*}$. Then secondly, $O(n^2m)$ many where for $u<v$ not adjacent we have $x_{uv}=\begin{cases*}$1$ & $u$ represents $v$\\ 0 & else\end{cases*}$. We want $k$ many color classes each having one representative.We also assume an ordering on the vertices and only allow the smallest vertex in each color class to be the representative. With these two, we can write a set of equations to describe the coloring problem where $\overline{N}(v)$ is the set of vertices that are non-adjacent to $v$.
	\end{frame}
	\begin{frame}
	\frametitle{Representative formulation}
	$x_u=\begin{cases*}$1$ & if $v$ is a representative\\ 0 & else\end{cases*}$\\
	$x_{uv}=\begin{cases*}$1$ & $u$ represents $v$\\ 0 & else\end{cases*}$
		\begin{alignat*}{2}
			&\sum_{u\in \overline{N}(v)}x_{uv}+x_v = 1 &\forall v\in V\\
			&x_{v,i}-x_{v,i+1}\geq 0 &\forall v\in V,\forall i\in\{1,...,k-1\}\\
			&x_{v,1}+x_{w,1}\geq 1 &\forall \{v,w\}\in E\\
			&x_{v,i-1}-x_{v,i}+x_{w,i-1}-x_{i}\leq 1 &\forall \{v,w\}\in E,\forall i\in \{2,...,k\}
		\end{alignat*}
	\end{frame}
	\begin{frame}
		\frametitle{Representative formulation}
		As logic formula this corresponds to the conjunction of the following clauses
		
		\begin{alignat*}{2}
			&\neg x_{v,k} &\forall v\in V\\
			&x_{v,i}\vee\neg x_{v,i+1} &\forall v\in V,\forall i\in\{1,...,k-1\}\\
			&x_{v,1}\vee x_{w,1} &\forall \{v,w\}\in E\\
			&\neg x_{v,i-1}\vee x_{v,i}\vee\neg x_{w,i-1}\vee x_{i} &\forall \{v,w\}\in E,\forall i\in \{2,...,k\}
		\end{alignat*}\\
		To forbid a cycle $C=(c_1,...c_l)\subseteq V$, we then add the clauses
		
		\begin{alignat*}{2}
			&\bigvee_{s=1,3,5,..l-1}\neg x_{c_s,i-1}\vee x_{c_s,i}\vee\bigvee_{s=2,4,6,..l}\neg x_{c_s,j-1}\vee x_{c_s,j} &\forall i\neq j\in\{1,...,k\}
		\end{alignat*}
		though, if one color w.l.o.g $i$ is $1$, there is no $x_{c_s,i-1}$ so leave that part out and just require $x_{c_s,i}$.
	\end{frame}
	\begin{frame}
		This leaves the problem of requiring, that at most $k$ many color classes i.e. representatives exist. For this we use the third set of variables $s_{j,i}=\begin{cases*}$1$ & $\{x_1,..x_j\}$ contains at least $i$ many $1$\\ 0 & else\end{cases*}$ and the following equations:
		\begin{alignat*}{2}
			&s_{1,1}=x_1 &\\
			&s_{1,i}=0 &\forall 2\leq i\leq k+1\\
			&s_{j-1,i}\leq s_{j,i} &\forall 2\leq j\leq n\forall 1\leq i\leq k+1\\
			&s_{j-1,i-1}+x_j\leq s_{j,i} &\forall 2\leq j\leq n\forall 1\leq i \leq k+1\\
			&s_{n,k+1}=0&\\
		\end{alignat*}\\
	\end{frame}
	\begin{frame}
		\frametitle{Runtime}
		Letting these run on Erdös-Rényi graphs with $n=20$ and varying density $d=0.3$ and $d=0.5$ we get the median:
		\begin{figure}[H]
			\centering
			\begin{minipage}{.5\textwidth}
				\centering
				\pgfplotstableread[col sep=semicolon]{Evaluating data 5.csv}\rawdata
				\resizebox{\columnwidth}{!}{
					\begin{tikzpicture}
						\begin{axis} [
							ybar,
							bar width=20pt,
							ylabel={Time (ms)},
							xlabel={density $0.3$},
							% Automatically read category text labels directly from the table
							xtick=data,
							xticklabels from table={\rawdata}{category},        xticklabel style={rotate=60, anchor=north east}],
							ymin=0, ymax=1500,
							enlarge x limits=0.3
							]
							
							% 2. Clean plot generation using structured table columns
							\addplot+[
							error bars/.cd,
							y dir=both,
							y explicit
							] table [
							x expr=\coordindex, 
							y=median2,
							% Dynamically calculate lower and upper error whisker margins
							y error minus expr=\thisrow{median2}-\thisrow{min},
							y error plus expr=\thisrow{max}-\thisrow{median2}
							] {\rawdata};
							
						\end{axis}
				\end{tikzpicture}}
			\end{minipage}%
			\begin{minipage}{.5\textwidth}
				\centering
				\pgfplotstableread[col sep=semicolon]{Evaluating data 6.csv}\rawdata
				\resizebox{\columnwidth}{!}{
					\begin{tikzpicture}
						\begin{axis} [
							ybar,
							bar width=20pt,
							ylabel={Time (ms)},
							xlabel={density $0.5$},
							% Automatically read category text labels directly from the table
							xtick=data,
							xticklabels from table={\rawdata}{category},        xticklabel style={rotate=60, anchor=north east}],
							ymin=0, ymax=250000,
							enlarge x limits=0.3
							]
							
							% 2. Clean plot generation using structured table columns
							\addplot+[
							error bars/.cd,
							y dir=both,
							y explicit
							] table [
							x expr=\coordindex, 
							y=median2,
							% Dynamically calculate lower and upper error whisker margins
							y error minus expr=\thisrow{median2}-\thisrow{min},
							y error plus expr=\thisrow{max}-\thisrow{median2}
							] {\rawdata};
							
						\end{axis}
				\end{tikzpicture}}
			\end{minipage}
		\end{figure}
	\end{frame}
	\begin{frame}
		\frametitle{Runtime}
		Letting all of these run on Erdös-Rényi graphs with $n=20$ and varying density $d=0.3$ and $d=0.5$ we get the median:
		\begin{figure}[H]
			\centering
			\begin{minipage}{.5\textwidth}
				\centering
				\pgfplotstableread[col sep=semicolon]{Evaluating data 5.csv}\rawdata
				\resizebox{\columnwidth}{!}{
					\begin{tikzpicture}
						\begin{axis} [
							ybar,
							bar width=20pt,
							ylabel={Time (ms)},
							xlabel={density $0.3$},
							% Automatically read category text labels directly from the table
							xtick=data,
							xticklabels from table={\rawdata}{category},        xticklabel style={rotate=60, anchor=north east}],
							ymin=0, ymax=40,
							enlarge x limits=0.3
							]
							
							% 2. Clean plot generation using structured table columns
							\addplot+[
							error bars/.cd,
							y dir=both,
							y explicit
							] table [
							x expr=\coordindex, 
							y=median2]{\rawdata};
							
						\end{axis}
				\end{tikzpicture}}
			\end{minipage}%
			\begin{minipage}{.5\textwidth}
				\centering
				\pgfplotstableread[col sep=semicolon]{Evaluating data 6.csv}\rawdata
				\resizebox{\columnwidth}{!}{
					\begin{tikzpicture}
						\begin{axis} [
							ybar,
							bar width=20pt,
							ylabel={Time (ms)},
							xlabel={density $0.5$},
							% Automatically read category text labels directly from the table
							xtick=data,
							xticklabels from table={\rawdata}{category},        xticklabel style={rotate=60, anchor=north east}],
							ymin=0, ymax=20000,
							enlarge x limits=0.3
							]
							
							% 2. Clean plot generation using structured table columns
							\addplot+[
							error bars/.cd,
							y dir=both,
							y explicit
							] table [
							x expr=\coordindex, 
							y=median2] {\rawdata};
							
						\end{axis}
				\end{tikzpicture}}
			\end{minipage}
		\end{figure}
	\end{frame}
	\begin{frame}
		\frametitle{Runtime}
		The runtime for thie last verison on Graphs with $n=40$ and density $0,5$:
		\pgfplotstableread[col sep=semicolon]{Evaluating data 7.csv}\rawdata
			\begin{tikzpicture}
				\begin{axis} [
					ybar,
					bar width=20pt,
					ylabel={Time (s)},
					xlabel={density $0.5$},
					% Automatically read category text labels directly from the table
					xtick=data,
					xticklabels from table={\rawdata}{category},
					ymin=0, ymax=25000,
					enlarge x limits=0.3
					]
					
					% 2. Clean plot generation using structured table columns
					\addplot+[
					error bars/.cd,
					y dir=both,
					y explicit
					] table [
					x expr=\coordindex, 
					y=median2,
					% Dynamically calculate lower and upper error whisker margins
					y error minus expr=\thisrow{median2}-\thisrow{min},
					y error plus expr=\thisrow{max}-\thisrow{median2}
					] {\rawdata};
					
				\end{axis}
			\end{tikzpicture}
	\end{frame}
	\begin{frame}{References} 
		\nocite{*}
		\bibliographystyle{plain}
		\bibliography{Literatur} 
	\end{frame}
\end{document}